\begingroup
\setlength{\LTcapwidth}{\textwidth}
\setlength{\tabcolsep}{4pt}
\renewcommand{\arraystretch}{1.12}
\begin{longtable}{@{}p{1.05cm}p{4.65cm}cccccc@{}}
\caption{Penilaian alternatif pendekatan menggunakan \textit{Weighted Sum Model}.}
\label{tab:tab3.5_penilaian_wsm} \\
\toprule
\textbf{Kode} & \textbf{Alternatif pendekatan} & \textbf{C1} & \textbf{C2} & \textbf{C3} & \textbf{C4} & \textbf{C5} & \textbf{Skor} \\
\midrule
\endfirsthead

\multicolumn{8}{@{}l}{\small\emph{Tabel \thetable{} (lanjutan)}} \\
\toprule
\textbf{Kode} & \textbf{Alternatif pendekatan} & \textbf{C1} & \textbf{C2} & \textbf{C3} & \textbf{C4} & \textbf{C5} & \textbf{Skor} \\
\midrule
\endhead

\bottomrule
\endlastfoot

P1 & Pemantauan manual berbasis ambang & 1 & 1 & 3 & 5 & 3 & 2,40 \\

P2 & Model prediksi tanpa penelusuran & 5 & 1 & 1 & 4 & 2 & 2,70 \\

P3 & Model bahasa besar tanpa penelusuran & 1 & 2 & 2 & 3 & 1 & 1,85 \\

P4 & RAG dengan kueri statis & 1 & 2 & 5 & 3 & 4 & 2,75 \\

P5 & Model fisiologis mekanistik & 5 & 3 & 3 & 1 & 3 & 3,10 \\

\textbf{P6} & \textbf{RAG dengan kueri terkondisi-prediksi} & \textbf{5} & \textbf{5} & \textbf{5} & \textbf{3} & \textbf{4} & \textbf{4,50} \\
\end{longtable}
\endgroup

\noindent 
    \textit{Keterangan: skor berada pada skala 1--5; 1 menunjukkan kesesuaian sangat rendah dan 5 menunjukkan kesesuaian sangat tinggi.}
